\documentclass{article}
\usepackage[left=2cm, right=2cm, top=2cm, bottom=2cm]{geometry}
\usepackage{graphicx}
\usepackage{amsmath}
\usepackage{amssymb}
\usepackage{amsthm}
\usepackage{fancyhdr}
\usepackage{verbatim}
\usepackage{listings}
\usepackage{xcolor}
\usepackage{pdfpages}
\usepackage{cancel}
\usepackage{physics}
\usepackage{esint}
\usepackage{braket}

\lstset{
    language=C++,
    basicstyle=\ttfamily\footnotesize,
    keywordstyle=\color{blue},
    commentstyle=\color{green},
    stringstyle=\color{red},
    numbers=left,
    numberstyle=\tiny\color{gray},
    frame=single,
    breaklines=true,
    captionpos=b,
}


\title{MTH 446: Lecture \# 16}
\author{Cliff Sun}

\newtheorem{theorem}{Theorem}[section]
\newtheorem{lemma}[theorem]{Lemma}
\newtheorem{definition}[theorem]{Definition}
\newtheorem{conjecture}[theorem]{Conjecture}
\newtheorem{proposition}[theorem]{Proposition}
\newtheorem{corollary}[theorem]{Corollary}
\newtheorem{one minute paper}[theorem]{One Minute Paper}

\pagestyle{fancy}
\lhead{\textbf{\thepage}\ \ \nouppercase{\rightmark}}
\chead{MTH 446: Lecture \# 16}
\rhead{Cliff Sun}

\begin{document}

\maketitle

\section*{Lecture Span}
\begin{enumerate}
    \item Complex Logs
\end{enumerate}

\section*{Complex logs}

Find $2\log(-1 + i)$. Note that 

\begin{align}
    -1 + i &= \sqrt{2}\left( -\frac{\sqrt{2}}{2} + i\frac{\sqrt{2}}{2} \right) \\
    &= \sqrt{2}\left( \cos(3\pi/4) + i\sin(3\pi/4) \right) \\
    &= \sqrt{2}\exp(i3\pi/4)
\end{align}

Therefore, 

\begin{equation}
    \arg(-1 + i) = 3\pi/4
\end{equation}

Therefore,

\begin{equation}
    \log(-1 + i) = \ln(\sqrt{2}) + i\frac{3\pi}{4}
\end{equation}

Thus, 

\begin{equation}
    2\log(-1 + i) = \ln(2) + i\frac{3\pi}{2}
\end{equation}

Note, $\log(z^2) = 2\log(z)$. 

\section*{Analyticity of the single-valued branch of $\log(z)$}

Consider the single-valued branch of $\log(z)$, defined as 

\begin{equation}
    \log(z) = \ln(r) + i\theta, \quad \alpha < \theta < 2\pi + \alpha
\end{equation}

Then, the domain of this log is $\mathbb{C}/\text{branch cut}$. Thus, $\log: D_{\alpha} \to \mathbb{C}$. Then, 

\begin{equation}
    \log(z) = u(r, \theta) + iv(r, \theta)
\end{equation}

Moreover, $u = \ln(r)$, and $v(r,\theta) = \theta$. Thus, 

\begin{align}
    u_r &= 1/r \\
    u_{\theta}  &= 0 \\
    v_r &= 0 \\
    v_{\theta} &= 1
\end{align}

The CR equations in Polar are 

\begin{equation}
    u_r = \frac{1}{r}v_{\theta}, \quad v_r = -\frac{1}{r}u_{\theta}
\end{equation}

The first equation is satisified. The second equation is also satisifed. Therefore, $\log(z)$ is complex differentiable except for the branch cut. 

\section*{Derivative of $\log(z)$}

Let $\log(z) = \ln r + i\theta$. Then recall that

\begin{equation}
    f'(z) = -\frac{i}{z}\left( u_{\theta} + iv_{\theta} \right)
\end{equation}

Therefore, 

\begin{equation}
    f'(z) = -\frac{i}{z}\left( 0 + i \right) = \frac{1}{z}
\end{equation}

If $z=0$ is called the branch point of $\log(z)$ (the point at which all the branch cuts intersect). 

\section*{Bernoulli's Sophism}

Bernoulii showed that $\log(z) = \log(-z)$. But this isn't a proof, so its a sophism. 

\section*{Identities}

Define, 

\begin{equation}
    c\log(z) = c\left[ \ln r + i(\theta + 2\pi k) \right]
\end{equation}

We can prove the following:
\begin{enumerate}
    \item $z^n = \exp(n \log z)$
    \item $z^{1/n} = \exp(\frac{1}{n}\log z)$
\end{enumerate}

\section*{The power function}

Note that $z^n = \exp(n\log(z))$. This motivates the following definition: 

\begin{equation}
    z^c = \exp(c\log z), \quad c \in \mathbb{R}
\end{equation}

Therefore, $z^c$ is a multiple valued function. The principle value is 

\begin{equation}
    P.V z^c = \exp(c\rm Log(z))
\end{equation}

Example, 

\begin{equation}
    i^i = \exp(i\log(i))
\end{equation}

Note that $i = \exp(i\pi/2)$. Then, 

\begin{equation}
    \exp(i(\ln(1) + i(\pi/2 + 2\pi k))) = \exp(-\frac{4k + 1}{2}\pi), k \in \mathbb{Z}
\end{equation}

Thus, 

\begin{equation}
    i^i = \exp(i\rm Log(i)) = \exp(-\pi/2)
\end{equation}

\subsection*{Example}

Let $\alpha = 3\pi/4$. Then, use the single valued log to evaluate $z^i$. Then, 

\begin{equation}
    i^i = \exp(i\log(i)) = \exp(i\theta), \quad \theta \in (\alpha, 2\pi + \alpha)
\end{equation}

Therefore, 

\begin{equation}
    \theta = 5\pi/2
\end{equation}

is the value we want. 

Evaluate 

\begin{equation}
    (1-i)^{4i}
\end{equation}

This is the same as:

\begin{equation}
    \exp(4i\log(1-i))
\end{equation}



\end{document}